\subsection{Torus Generalization (PhaseDial)}
\label{sec:torus_generalization}

\paragraph{Experiment Description.}
The Torus Generalization experiment evaluates whether frequency-domain
splits of the PhaseDial fingerprint produce \emph{super-additive}
analogical gain---that is, whether a two-axis torus rotation applied to
a composite fingerprint $\phi_A \star \phi_B$ can locate a third concept
$C$ that is simultaneously close to both $A$ and $B$, exceeding the
single-axis ceiling achievable by rotating $\phi_A$ or $\phi_B$ alone.

Three split strategies are compared:
\begin{itemize}[nosep]
  \item $T^2$-\textbf{contiguous}: partition the frequency spectrum at a
        fixed index.
  \item $T^2$-\textbf{random}: randomly assign frequency bins to each
        axis (fixed seed).
  \item $T^2$-\textbf{energy}: assign bins to axes greedily by the
        energy difference $|E_A(k) - E_B(k)|$, placing the
        most discriminative frequencies on separate axes.
\end{itemize}

The experiment sweeps all grid angles $(\alpha, \beta) \in [0, 2\pi)^2$
at 5\textdegree{} resolution and records the best gain $g = \min(
\mathrm{sim}(\phi_C, \phi_A), \mathrm{sim}(\phi_C, \phi_B))$ for each
configuration.

\paragraph{Status.}
The torus sweep computation has been temporarily suspended pending the
stabilisation of the \texttt{geometry.HybridSubstrate} pointer-based
query API introduced in the current iteration.  The experiment infrastructure
(split strategies, grid sweep, super-additivity test) is fully implemented
and will resume automatically once the API is finalised.
The section will be populated with results and figures on the next
experimental run.
